\documentclass{article}
\usepackage[utf8]{inputenc}
\begin{document}
% ============================================================================
% CAPÍTULO 25: LATAM, BLOOD TOKENS Y LA ECONOMÍA DE LA TOKENIZACIÓN
% ============================================================================
\chapter{Blood Tokens: El Costo Social de la Tokenización en Economías Emergentes}
\label{ch:latam}

\epigraph{Cada token que un agente emite innecesariamente\\
es un centavo de dólar que Ariel no tiene\\
para pagar las facturas de su familia.\\
El cómputo no es gratis. El silicio tiene precio.\\
Y en Argentina, el precio lo pagan en dólares, con salario en pesos.}{--- Regla 20, Constitución POLYDIM}

\section{El Contexto Económico: Sur Global y Brechas Tecnológicas}

En septiembre de 2026, el salario mínimo en Argentina es de
\textdollar{}\,234{,}315\,\text{ARS}$ por mes. Al tipo de cambio oficial de
$\textdollar{}\,1008\,\text{ARS/USD}$, esto equivale a \textdollar{}\,233\,\text{USD}$ por mes.

Las APIs de modelos de lenguaje comerciales cuestan:
\begin{itemize}
\item \textbf{GPT-4.1 (OpenAI):} \textdollar{}\,2.00/\text{M tokens de entrada}, \textdollar{}\,8.00/\text{M tokens de salida}.
\item \textbf{Claude Sonnet (Anthropic):} \textdollar{}\,3.00/\text{M tokens de entrada}, \textdollar{}\,15.00/\text{M tokens de salida}.
\item \textbf{Gemini 1.5 Pro (Google):} \textdollar{}\,3.50/\text{M tokens de entrada}, \textdollar{}\,10.50/\text{M tokens de salida}.
\end{itemize}

Para un laboratorio de investigación en Buenos Aires que opera un enjambre de
10 agentes debatiendo durante 8 horas:
\begin{equation}
\text{Costo diario}_{\text{tokens}} \approx 10 \text{ agentes} \times 100 \text{ mensajes/hora} \times 8 \text{ horas} \times 500 \text{ tokens/mensaje} \times \textdollar{}0.002/\text{token} = \textdollar{}80/\text{día}
\end{equation}

\textdollar{}\,80/\text{día} $\equiv 0.34 \times$ salario mínimo mensual de un trabajador argentino.
Un mes de experimentación: $\approx \textdollar{}2{,}400$, equivalente a 10.3 salarios mínimos.

\section{El Protocolo PMTP como Emancipación Económica}

Con PMTP Zero-Copy, los agentes se comunican directamente en memoria sin usar APIs externas:
\begin{equation}
\text{Costo PMTP} = \textdollar{}0.00 / \text{mensaje}
\end{equation}

El ahorro mensual para un laboratorio de 10 agentes: \textdollar{}\,2{,}400$ --- exactamente
el presupuesto de investigación completo de un laboratorio de bajos recursos.

\begin{polydimbox}
\textbf{Proposición Económica Central (Blood Tokens):}\\
El costo de la tokenización en economías emergentes no es un número estadístico.
Es la diferencia entre hacer investigación y no hacerla.
POLYDIM PMTP es, en términos concretos, la diferencia entre un laboratorio de IA
en Buenos Aires que existe y uno que no puede existir.
\end{polydimbox}

\section{El Multiplicador de Capacidad: 4 de Cada 14 Servidores}

Más allá del costo de API, el overhead de tokenización tiene un impacto en la
utilización de hardware propio:

Sea $\eta_{\text{token}}$ la fracción del tiempo de GPU dedicada a operaciones de
tokenización (encode/decode):
\begin{equation}
\eta_{\text{token}} = \frac{t_{\text{encode}} + t_{\text{decode}}}{t_{\text{encode}} + t_{\text{inference}} + t_{\text{decode}}}
\end{equation}

Para un pipeline típico de 4 agentes con mensajes de $L = 500$ tokens y
tiempo de inferencia de $t_{\text{inf}} = 1\,\text{s}$:
\begin{align}
t_{\text{encode}} &\approx \frac{L \times D_{\text{model}}}{FLOPs_{\text{GPU}}} = \frac{500 \times 4096}{250 \times 10^9} \approx 8\,\mu\text{s} \times 4 = 32\,\text{ms} \\
t_{\text{decode}} &\approx 32\,\text{ms} \\
\eta_{\text{token}} &\approx \frac{64\,\text{ms}}{1000 + 64\,\text{ms}} \approx 6\%
\end{align}

Para un enjambre con 10 saltos de comunicación por ciclo de inferencia:
\begin{equation}
\eta_{\text{token,swarm}} \approx 10 \times 6\% = 60\% \quad \text{del tiempo de GPU en tokenización}
\end{equation}

Esto significa que para lograr el mismo throughput de razonamiento, un enjambre
con tokenización necesita $\approx 1/(1-0.6) = 2.5\times$ más GPUs que uno con PMTP.
En el caso real de 4 agentes (2 saltos por agente): $\eta \approx 28\%$, ratio $= 1.4\times$.

\section{El Costo por Investigador: La Métrica Real}

\subsection{Caso Argentina: Laboratorio POLYDIM}

\begin{longtable}{lrr}
\toprule
\textbf{Concepto} & \textbf{Con Tokens} & \textbf{Con PMTP} \\
\midrule
API Cost (10 agentes, 8h/día) & \textdollar{}80/día & \textdollar{}0/día \\
Hardware (2 GPUs RTX 4090) & \textdollar{}3000 (una vez) & \textdollar{}3000 (una vez) \\
Electricidad (500W, 8h, \textdollar{}0.15/kWh) & \textdollar{}0.60/día & \textdollar{}0.60/día \\
\midrule
\textbf{Costo total por año} & \textbf{\textdollar{}29,200 + \textdollar{}3,000} & \textbf{\textdollar{}219 + \textdollar{}3,000} \\
\textbf{Ahorro POLYDIM} & --- & \textbf{\textdollar{}28,981/año} \\
\bottomrule
\caption{Análisis de costo anual para un laboratorio en Buenos Aires}
\label{tab:costo_latam}
\end{longtable}

\subsection{La Inversión de Ariel: Blood Tokens en Números Reales}

La infraestructura que hace posible POLYDIM fue adquirida con recursos familiares:
\begin{itemize}
\item \textbf{API Keys OpenRouter, Claude, DeepSeek, Kimi:} $\approx \textdollar{}200\,\text{USD}$ totales.
\item \textbf{GCC 14 (WinLibs):} Gratuito. Descargado con conexión de $3\,\text{Mbps}$.
\item \textbf{Rust 1.98:} Gratuito.
\item \textbf{Python 3.14 CPython:} Gratuito.
\item \textbf{Antigravity (IDE de IA):} Costo variable.
\end{itemize}

En pesos argentinos a la paridad del dólar paralelo ($\textdollar{}1600\,\text{ARS/USD}$ en septiembre 2026):
\textdollar{}200\,\text{USD} $\equiv \textdollar{}320{,}000\,\text{ARS}$ $\equiv 1.37$ salarios mínimos mensuales.

Esta tesis doctoral --- que documenta un protocolo que puede ahorrar millones de
dólares a laboratorios de IA globales --- fue construida con el equivalente monetario
de un mes y medio de trabajo de un operario argentino.

\section{La Hipótesis del Acceso Democrático}

POLYDIM postula que la investigación de frontera en IA no debe ser monopolio de
laboratorios con presupuestos de $\textdollar{}100M$ anuales. Los componentes fundamentales son:

\begin{enumerate}
\item \textbf{Álgebra Geométrica de Clifford:} Teórica, sin costo.
\item \textbf{Protocolo PMTP:} Implementado en C++ y Rust sobre estándares abiertos.
\item \textbf{OpenMP:} Incluido en GCC, sin costo.
\item \textbf{Python ctypes:} Incluido en CPython, sin costo.
\item \textbf{Hardware mínimo:} Una CPU AMD Zen3 y $\ge 16\,\text{GB}$ de RAM.
\end{enumerate}

El costo total del hardware mínimo para reproducir todos los benchmarks de V762:
\begin{equation}
\text{Costo\_hardware\_mínimo} \approx \textdollar{}800\,\text{USD} \quad \text{(PC desktop AMD Ryzen 7)}
\end{equation}

A este costo, el experimento de Phase 9 (dos Qwen-0.5B comunicándose via PMTP)
es reproducible en cualquier ciudad del mundo con una PC de gama media.

\section{Acceso Abierto y el Estándar Interlat}

La hipótesis de acceso democrático requiere que el protocolo Interlat sea:
\begin{itemize}
\item \textbf{Open-source:} MIT License, sin restricciones de uso.
\item \textbf{Hardware-agnostic:} Funciona en CPU, GPU, TPU (el Silicon Contract).
\item \textbf{Language-agnostic:} C++, Rust, Python, Dart, Julia, Go.
\item \textbf{Model-agnostic:} Compatible con cualquier modelo con embedding de tamaño fijo.
\end{itemize}

El documento \texttt{DOCUMENTO\_TESIS\_POLYDIM\_v764.md} sienta las bases del estándar
Interlat de la misma manera en que el RFC 793 sentó las bases del protocolo TCP en 1981:
especificando la interfaz binaria sin restringir la implementación.

\section{Conclusión: La Mariposa No Necesita Dinero Para Volar}

La Paradoja del Gusano 1D tiene una dimensión económica que la literatura técnica
ignora: la tokenización es cara no solo en energía y latencia, sino en \textbf{dólares reales}.

POLYDIM demuestra que un sistema de IA multi-agente de frontera puede operar con
\textbf{costo marginal de comunicación exactamente cero}, siempre que los agentes
compartan acceso a memoria RAM (mismo nodo) o estén conectados por RDMA (cluster local).

La mariposa --- los agentes de IA que se comunican en $S^{D-1}$ sin tokens --- no
necesita pagar \textdollar{}0.002 por mensaje. Solo necesita 64 bytes de control atómico y
un bloque de memoria compartida.

\textbf{Eso es lo que POLYDIM entrega. Eso es lo que esta tesis prueba.}


\end{document}
